\documentclass[10pt,a4paper]{amsart}
\usepackage[margin=19mm]{geometry}
\usepackage{amsmath,amssymb,mathtools}
\usepackage[scr=rsfs,cal=euler]{mathalfa}
\usepackage{xfrac}
\usepackage[unicode,hidelinks,bookmarks=false]{hyperref}
\newcommand{\ie}{i.e.\ }
\newcommand{\cf}{cf.\ }
\newcommand{\resp}{respectively\ }
\newcommand{\Subsection}{\subsection}
\providecommand{\nobreakdash}{\nobreak\hbox{-}}
\setlength{\emergencystretch}{2em}
\multlinegap0pt
\usepackage{bm}
\usepackage{bbm}
\usepackage{dsfont}
\usepackage{xspace}
\usepackage{cancel}
\usepackage{mathtools}
\usepackage{amsthm}
\makeatletter
\@ifundefined{Assum}{\newtheorem{Assum}{Assumption}}{}
\@ifundefined{Lem}{\newtheorem{Lem}{Lemma}}{}
\@ifundefined{Thm}{\newtheorem{Thm}{Theorem}}{}
\makeatother
\DeclareMathOperator{\divg}{div}
\title[Two improved Liouville-type theorems for the 3D stationary t]{Two improved Liouville-type theorems for the 3D stationary tropical climate model without temperature assumptions}
\author{Zhibing Zhang}
\date{Source: arXiv:2608.13149v1 (CC BY 4.0)}
\begin{document}
\maketitle

\noindent\emph{Source and licence.} Two improved Liouville-type theorems for the 3D stationary tropical climate model without temperature assumptions. Version arXiv:2608.13149v1 (CC BY 4.0), distributed under the Creative Commons Attribution 4.0 International licence, as recorded in the arXiv licence metadata for this exact version and shown on the abstract page https://arxiv.org/abs/2608.13149v1. Full rights notice: licenses/arxiv-2608.13149-CC-BY-4.0.txt.\par\smallskip

\noindent\textit{Editorial scope.} Counted unit: the Lem whose source label is \texttt{Lem4.9} in the article (source0.tex lines 1339--1346, source label \texttt{Lem4.9}), delivered together with the article's own complete original proof of it (source0.tex lines 1347--1540, one \texttt{proof} environment of 194 lines). No mathematics is written, repaired or re-derived: statement and proof are the article's own text line for line. Required context retained uncounted: equ1.1 (lines 132--141); a1.1 (lines 211--220); a1.2 (lines 223--235); ine1.3 (lines 226--228); a1.3 (lines 238--240); a1.4 (lines 242--259); main2 (lines 272--280); Lem2.5 (lines 367--373); ine4.2 (lines 988--990); ine4.8 (lines 1041--1044); Lem4.3 (lines 1085--1090); ine4.14 (lines 1087--1089); ine4.15 (lines 1093--1095); ine4.16 (lines 1110--1112); ine4.17 (lines 1128--1130); Lem4.6 (lines 1158--1171); ine4.19 (lines 1174--1178); ine4.30 (lines 1285--1287). Every definition, display and label of those lines is retained unchanged. Omitted from this package and not redistributed: 1--131, 142--210, 221--222, 229--237, 241--241, 260--271, 281--366, 374--987, 991--1040, 1045--1084, 1090--1092, 1096--1109, 1113--1127, 1131--1157, 1172--1173, 1179--1284, 1288--1338, 1541--1678. Nothing inside a retained excerpt is omitted or altered: every retained body is delivered exactly as the article's lines are written, so no omission receipt of any kind accompanies this package. Citations: the retained excerpts carry no citation command at all, so no bibliography entry of the article is transcribed and no reference is redistributed. Reference adaptation: none was needed and none was made; every reference inside the delivered unit resolves inside it, so no unresolved reference or citation is produced. Printed slips kept literal: no printed word, symbol, hypothesis or number of the counted statement or of its proof was altered. Layout and class: the article's own class is \texttt{amsart} with packages including mathrsfs, cite, amsmath, amssymb, amsfonts, color; the wrapper is the lane's pinned \texttt{amsart} editorial wrapper at 10pt on A4, which restates the theorem-like environments the retained text needs and adds the article's own macro definitions that the retained text needs (\texttt{\textbackslash divg}), with extra wrapper packages (bm, bbm, dsfont, xspace, cancel, mathtools). The delivered numbering is the unit's own. Assets: no retained span contains a figure environment, a graphics-inclusion command, a PDF-page inclusion, a table or a TikZ picture; no image file is copied into this package, the package declares no asset dependency and no image file is redistributed with it. Licence: version arXiv:2608.13149v1 of this article is distributed under the Creative Commons Attribution 4.0 International licence, as recorded in the arXiv licence metadata for this exact version and shown on the abstract page https://arxiv.org/abs/2608.13149v1. A full rights notice is delivered at licenses/arxiv-2608.13149-CC-BY-4.0.txt. Characters: every retained excerpt is pure ASCII, so the delivered engine, XeLaTeX with its default Latin Modern text font, reports no missing character.
\begin{equation}\label{equ1.1}
  \left\{
    \begin{array}{ll}
     -\Delta u+(u\cdot\nabla) u+\nabla \pi+\divg(v\otimes v)=0,  \\
   -\Delta v+(u\cdot\nabla) v+\nabla \theta+(v\cdot\nabla)u=0, \\
   -\Delta \theta+u\cdot\nabla\theta+\divg v=0,\\
     \divg u=0,
\end{array}
  \right.
\end{equation}

\begin{Assum}\label{a1.1}
Let $(p,q)\in\left(\frac{3}{2},3\right]\times[1,2]$ and $\alpha\in\left[0,\frac{2}{p}-\frac{1}{3}\right]$, $\beta\in\left[0,\frac{3}{q}-\frac{1}{2}\right]$, $\gamma\in\left[0,\gamma_p\right)$, where
\begin{equation*}
\gamma_p=
\begin{cases}
\frac{6-p}{18-6p}, & \text{ if }\;\frac{3}{2}<p<3, \\
+\infty, & \text{ if }\;p=3.
\end{cases}
\end{equation*}
\end{Assum}

\begin{Assum}\label{a1.2}
\begin{itemize}
\item[(i)] When $\frac{6-2p}{6-p} + \frac{6-3q}{6-q} < 1$, we require
\begin{align}\label{ine1.3}
\frac{p}{6-p}\alpha+\frac{2q}{6-q}\beta\leq1.
\end{align}
\item[(ii)] When $\frac{6-2p}{6-p} + \frac{6-3q}{6-q} \geq 1$ and $\gamma=0$, we require that \eqref{ine1.3} holds.
\item[(iii)] When $\frac{6-2p}{6-p} + \frac{6-3q}{6-q} \geq 1$ and $\gamma>0$, we require
 \begin{align}\label{ine1.6}
 \frac{p}{6-p}\alpha+\frac{2q}{6-q}\beta+\left(\frac{6-2p}{6-p}+\frac{6-3q}{6-q}-1\right)\gamma<1.
 \end{align}
\end{itemize}
\end{Assum}

\begin{Assum}\label{a1.3}
Let $\lambda\in\left[0,\frac{3}{p}-1\right]$ and $\mu\geq0$.
\end{Assum}

\begin{Assum}\label{a1.4}
\begin{itemize}
\item[(i)]
When $\frac{6-2p}{6-p} + \frac{6-3q}{6-q} < 1$ and $\frac{p}{6-p}\alpha+\frac{2q}{6-q}\beta=1$, we require
$$
\frac{p}{6-p}\lambda+\frac{2q}{6-q}\mu\leq \frac{3-p}{6-p}+\frac{6-3q}{12-2q}+\frac{1}{2}.
$$	
\item[(ii)] When $\frac{6-2p}{6-p} + \frac{6-3q}{6-q} \geq 1$ and $\gamma=0$,
we require
$$
\frac{p}{6-p}\lambda+\frac{2q}{6-q}\mu<1.
$$	
\item[(iii)] When $\beta=\frac{3}{q}-\frac{1}{2}$, we require
$$
p\neq3\;\;\text{ and }\;\;\frac{p}{6-p}\lambda+2\mu\leq \frac{3-p}{6-p}.
$$
\end{itemize}	
\end{Assum}

\begin{Thm}\label{main2}
Let $(u,\pi,v,\theta)$ be a smooth solution of \eqref{equ1.1}. Suppose the parameters $p$, $q$, $\alpha$, $\beta$, $\gamma$, $\lambda$, $\mu$ satisfy Assumptions \ref{a1.1}, \ref{a1.2}, \ref{a1.3} and \ref{a1.4}.
Furthermore, assume that one of the following conditions holds
\begin{align*}
&\mathrm{(B1)}\;\limsup\limits_{R\rightarrow+\infty}\left[X_{p,\alpha,\lambda}(R)+Y_{q,\beta,\mu}(R)+Y'_{2,\gamma}(R)\right]<+\infty,\text{ with }p\in\left(\frac{3}{2},3\right),\\
&\mathrm{(B2)}\;\lim\limits_{R\rightarrow+\infty}X_{p,\alpha}(R)=0,\;\limsup\limits_{R\rightarrow+\infty}\left[Y_{q,\beta,\mu}(R)+Y'_{2,\gamma}(R)\right]<+\infty,\text{ with }p=3.
\end{align*}
Then $u=v=0$ and $\theta$ is a constant.
\end{Thm}

\begin{Lem}\label{Lem2.5}
Let $(u,\pi,v,\theta)$ be a smooth solution of \eqref{equ1.1}. For each $L>1$, $q\geq1$ and $\rho>0$, there exists a positive constant $C_L$  such that
\begin{align}\label{ine2.1}
\|\theta - (\theta)_{E}\|_{L^2(E)} \leq C_L \|u\|_{L^3(E)} \left( \|\nabla v\|_{L^2(E)} + \rho^{\frac{1}{2}-\frac{3}{q}} \|v\|_{L^q(E)} \right) + C_L\|\nabla v\|_{L^2(E)},
\end{align}
where $E= B_{L\rho}\backslash B_\rho$.
\end{Lem}

	\begin{align}\label{ine4.2}
		E'(R)\geq\frac{1}{R}\int_{A_R}\left(|\nabla u|^{2}+|\nabla v|^{2}+|\nabla \theta|^{2}\right)dx.
	\end{align}

	\begin{align}
		\|u\|_{L^2(A_R)}&\leq CR\|\nabla u\|_{L^2 (A_R)}+CR^{\frac{3}{2}-\frac{3}{p}}\|u\|_{L^p(A_R)},\label{ine4.7}\\
        \|v\|_{L^2(A_R)}&\leq CR\|\nabla v\|_{L^2 (A_R)}+CR^{\frac{3}{2}-\frac{3}{q}}\|v\|_{L^q(A_R)}.\label{ine4.8}
	\end{align}

\begin{Lem}\label{Lem4.3}
Let the assumptions be the same as those in {\rm Theorem \ref{main2}}. Then there exist three positive constants $R_1>3$, $\tau_5$ and $C$ such that
	\begin{align}\label{ine4.14}
		f(2R)\leq C(\ln R)^{\tau_5},\;\forall R\geq R_1.
	\end{align}
\end{Lem}

\begin{align}
&\|u\|_{L^p\left(A_R\right)}\leq CR^\alpha(\ln R)^\lambda,\quad\|v\|_{L^q\left(A_R\right)}\leq CR^\beta(\ln R)^\mu,\quad\|\nabla v\|_{L^2\left(A_R\right)}\leq CR^\gamma.\label{ine4.15}
\end{align}

	\begin{align}\label{ine4.16}
	K_1\leq&C\left[R\ln R E'(R)\right]^\frac{\tau_5}{1+\tau_5},\;\forall R\geq R_1.
	\end{align}

\begin{equation}\label{ine4.17}
K_2+K_3\leq C[R\ln R E'(R)]^\frac{1}{2},\;\forall R\geq R_1.
\end{equation}

\begin{Lem}\label{Lem4.6}
Let the assumptions be the same as those in {\rm Theorem \ref{main2}}. Assume $E(R)\not\equiv0$. Then there exist two positive constants $R_5>3$ and $C$ such that
\begin{align}\label{ine4.18}
K_4\leq \frac{1}{16}E(R)+C\chi(p)\left[R\ln R E'(R)\right]^{\frac{9-3p}{6-p}},\;\forall R\geq R_5,
\end{align}
where $\chi(p)$ is defined by
\begin{equation*}
\chi(p)=
\begin{cases}
1, & \text{ if  }\frac{3}{2}<p<3, \\
0, & \text{ if  } p=3.
\end{cases}
\end{equation*}
\end{Lem}

	\begin{align}\label{ine4.19}
		\|u\|_{L^6(A_R)}&\leq \|u-(u)_{A_R}\|_{L^6(A_R)}+\|(u)_{A_R}\|_{L^6(A_R)}\notag\\
		&\leq C\|\nabla u\|_{L^2 (A_R)}+CR^\frac{1}{2}\left|(u)_{A_R}\right|\\
		&\leq C\|\nabla u\|_{L^2 (A_R)}+CR^{\frac{1}{2}-\frac{3}{p}}\|u\|_{L^p(A_R)}.\notag
	\end{align}

\begin{align}\label{ine4.30}
\|v\|_{L^6(A_R)}&\leq C\|\nabla v\|_{L^2 (A_R)}+CR^{\frac{1}{2}-\frac{3}{q}}\|v\|_{L^q(A_R)}.
\end{align}

\begin{Lem}\label{Lem4.9}
Let the assumptions be the same as those in {\rm Theorem \ref{main2}}. Let $\tau_5$ be  the constant given by  {\rm Lemma \ref{Lem4.3}}. Let $\chi(p)$ be the function defined in {\rm Lemma \ref{Lem4.6}}. Assume $E(R)\not\equiv0$. Then there exist three positive constants $\tau_9\in(0,1)$, $R_{11}>3$ and $C$ such that for any $R\geq R_{11}$, it holds that
\begin{align}\label{ine4.35}
K_7 \leq &\frac{1}{8} E(R)+C\left[R\ln RE'(R)\right]^{\tau_9}+C\left[R\ln RE'(R)\right]^{\frac{3-p}{6-p}+\frac{1}{2}}\notag\\
&+C\chi(p)\left(\left[R\ln RE'(R)\right]^{\frac{3-p}{6-p}+\frac{6-3q}{12-2q}}+\left[R\ln RE'(R)\right]^{\frac{3-p}{6-p}}\right)\\
&+C\left[R\ln RE'(R)\right]^{\frac{6-3q}{12-2q}+\frac{1}{2}}+C[R\ln RE'(R)]^\frac{\tau_5}{1+\tau_5}+C[R\ln RE'(R)]^\frac{1}{2}.\notag
\end{align}
\end{Lem}

\begin{proof}
By Lemma \ref{Lem2.5}, we have
\begin{align}\label{ine4.36}
		K_7 \leq& C R^{-1} \|v\|_{L^2(A_R)}\|u\|_{L^3(A_R)}\|\nabla v\|_{L^2(A_R)}+CR^{-\frac{3}{q}-\frac{1}{2}}\|v\|_{L^2(A_R)}\|u\|_{L^3(A_R)}\|v\|_{L^q(A_R)}\notag\\
        &+C R^{-1} \|v\|_{L^2(A_R)}\|\nabla v\|_{L^2(A_R)}\\
		=&:K_{71} + K_{72} + K_{73}.\notag
\end{align}


We first derive an estimate for $K_{71}$. When $p=3$, we deduce from  Assumptions \ref{a1.3} and \ref{a1.4}(iii) that $\lambda=0$ and $\beta<\frac{3}{q}-\frac{1}{2}$. Using Assumption \ref{a1.4}(i), we conclude
\begin{align*}
&R^{\alpha+\frac{2q}{6-q}\beta-1}(\ln R)^{\frac{2q}{6-q}\mu}\leq C(\ln R)^{\frac{6-3q}{12-2q}+\frac{1}{2}},\\
&R^{\alpha+\frac{2q}{6-q}\beta-1+\frac{6-3q}{6-q}\left[\beta-\left(\frac{3}{q}-\frac{1}{2}\right)\right]}(\ln R)^{\mu}\leq C(\ln R)^\frac{1}{2}.
\end{align*}
By the interpolation inequality, \eqref{ine4.30}, \eqref{ine4.15}, \eqref{ine4.2} and the above two inequalities, we obtain
\begin{align}\label{ine4.36a}
		K_{71} \leq& C R^{-1}\|u\|_{L^3(A_R)}\|v\|_{L^q(A_R)}^{\frac{2q}{6-q}}\|v\|_{L^6(A_R)}^{\frac{6-3q}{6-q}}\|\nabla v\|_{L^2(A_R)}\notag\\
               \leq& C R^{-1}\|u\|_{L^3(A_R)}\|v\|_{L^q(A_R)}^{\frac{2q}{6-q}}\left(\|\nabla v\|_{L^2(A_R)}^{\frac{6-3q}{6-q}}+\left(R^{\frac{1}{2}-\frac{3}{q}} \|v\|_{L^q(A_R)}\right)^{\frac{6-3q}{6-q}}\right)\|\nabla v\|_{L^2(A_R)}\notag\\
               \leq&C R^{\alpha+\frac{2q}{6-q}\beta-1}(\ln R)^{\frac{2q}{6-q}\mu}\left[R E'(R)\right]^{\frac{6-3q}{12-2q}+\frac{1}{2}}\\
               &+C R^{\alpha+\frac{2q}{6-q}\beta-1+\frac{6-3q}{6-q}\left[\beta-\left(\frac{3}{q}-\frac{1}{2}\right)\right]}(\ln R)^{\mu}\left[R E'(R)\right]^{\frac{1}{2}}\notag\\
               \leq&C\left[R\ln RE'(R)\right]^{\frac{6-3q}{12-2q}+\frac{1}{2}}+C\left[R\ln R E'(R)\right]^{\frac{1}{2}}.\notag
\end{align}


When $\frac{3}{2}<p<3$, the estimate for $K_{71}$ becomes more involved. By the interpolation inequality, \eqref{ine4.19} and \eqref{ine4.30}, we obtain
\begin{align}\label{ine4.37}
	K_{71} \leq& C R^{-1} \|u\|_{L^p(A_R)}^{\frac{p}{6-p}}\|v\|_{L^q(A_R)}^{\frac{2q}{6-q}} \|u\|_{L^6(A_R)}^{\frac{6-2p}{6-p}}\|v\|_{L^6(A_R)}^{\frac{6-3q}{6-q}}  \|\nabla v\|_{L^2(A_R)}\notag\\
	\leq &C R^{-1} \|u\|_{L^p(A_R)}^{\frac{p}{6-p}}\|v\|_{L^q(A_R)}^{\frac{2q}{6-q}}\|\nabla u\|_{L^2(A_R)}^{\frac{6-2p}{6-p}}\|\nabla v\|_{L^2(A_R)}^{\frac{6-3q}{6-q}}\|\nabla v\|_{L^2(A_R)}\notag\\
&+C R^{-1} \|u\|_{L^p(A_R)}^{\frac{p}{6-p}}\|v\|_{L^q(A_R)}^{\frac{2q}{6-q}}\|\nabla u\|_{L^2(A_R)}^{\frac{6-2p}{6-p}}\left(R^{\frac{1}{2}-\frac{3}{q}} \|v\|_{L^q(A_R)}\right)^{\frac{6-3q}{6-q}}\|\nabla v\|_{L^2(A_R)}\\
&+C R^{-1} \|u\|_{L^p(A_R)}^{\frac{p}{6-p}}\|v\|_{L^q(A_R)}^{\frac{2q}{6-q}}\left(R^{\frac{1}{2}-\frac{3}{p}} \|u\|_{L^p(A_R)}\right)^{\frac{6-2p}{6-p}}
\|v\|_{L^6(A_R)}^{\frac{6-3q}{6-q}}\|\nabla v\|_{L^2(A_R)}\notag\\
=&:K_{711} + K_{712} + K_{713}.\notag
\end{align}


To begin with, we establish a bound for $K_{711}$ via case-by-case analysis.
When $\frac{6-2p}{6-p} + \frac{6-3q}{6-q} < 1$, using \eqref{ine4.15}, \eqref{ine4.2} and Assumptions \ref{a1.2}(i), \ref{a1.4}(i), we obtain
\begin{align}\label{ine4.38}
	K_{711} \leq &C(R)^{\frac{p}{6-p}\alpha+\frac{2q}{6-q}\beta-1}(\ln R)^{\frac{p}{6-p}\lambda+\frac{2q}{6-q}\mu}\left[R E'(R)\right]^{\tau_6}\leq C\left[R \ln R E'(R)\right]^{\tau_6},
\end{align}
where $\tau_6$ is defined by
$$
\tau_6=\frac{3-p}{6-p}+\frac{6-3q}{12-2q}+\frac{1}{2}.
$$
When $\frac{6-2p}{6-p} + \frac{6-3q}{6-q}\geq1$ and $\gamma>0$, it follows from Assumption \ref{a1.2}(iii) that
$$
\frac{p}{6-p}\alpha+\frac{2q}{6-q}\beta<1.
$$
Using \eqref{ine4.15}, \eqref{ine4.14} and \eqref{ine4.2}, we get
\begin{align}\label{ine4.39}
K_{711} \leq&C R^{-1} \|u\|_{L^p(A_R)}^{\frac{p}{6-p}}\|v\|_{L^q(A_R)}^{\frac{2q}{6-q}}\left[f(2R)\right]^{\frac{3-p}{6-p}+\frac{6-3q}{12-2q}}\|\nabla v\|_{L^2(A_R)}\notag\\
\leq&C R^{\frac{p}{6-p}\alpha+\frac{2q}{6-q}\beta-1}(\ln R)^{\frac{p}{6-p}\lambda+\frac{2q}{6-q}\mu+\left(\frac{3-p}{6-p}+\frac{6-3q}{12-2q}\right)\tau_5} \left[RE'(R)\right]^{\frac{1}{2}}\notag\\
\leq&C\left[RE'(R)\right]^{\frac{1}{2}}\\
\leq&C\left[R\ln RE'(R)\right]^{\frac{1}{2}}.\notag
\end{align}
When $\frac{6-2p}{6-p} + \frac{6-3q}{6-q}\geq1$ and $\gamma=0$, thanks to Assumption \ref{a1.4}(ii) and the fact that $\frac{6-2p}{6-p} + \frac{6-3q}{6-q}\in[1,2)$, we can choose a positive number
$\tau_7$ such that
$$
2\left(\frac{p}{6-p}\lambda+\frac{2q}{6-q}\mu\right)-\frac{6-2p}{6-p}-\frac{6-3q}{6-q}<\tau_7<2-\frac{6-2p}{6-p} -\frac{6-3q}{6-q}.
$$
We denote
$$
\tau_8=\frac{\tau_7}{2}+\frac{1}{2}\left(\frac{6-2p}{6-p} + \frac{6-3q}{6-q}\right).
$$
It is not difficult to verify that
$$
\frac{p}{6-p}\lambda+\frac{2q}{6-q}\mu<\tau_8<1.
$$
Using \eqref{ine4.15} and \eqref{ine4.2}, we conclude that
\begin{align}\label{ine4.40}
K_{711} \leq&C R^{-1} \|u\|_{L^p(A_R)}^{\frac{p}{6-p}}\|v\|_{L^q(A_R)}^{\frac{2q}{6-q}}\|\nabla v\|_{L^2(A_R)}^{1-\tau_7}\|\nabla u\|_{L^2(A_R)}^{\frac{6-2p}{6-p}}\|\nabla v\|_{L^2(A_R)}^{\frac{6-3q}{6-q}}\|\nabla v\|_{L^2(A_R)}^{\tau_7}\notag\\
\leq&C R^{\frac{p}{6-p}\alpha+\frac{2q}{6-q}\beta-1}(\ln R)^{\frac{p}{6-p}\lambda+\frac{2q}{6-q}\mu} \left[RE'(R)\right]^{\tau_8}\\
\leq&C\left[R\ln RE'(R)\right]^{\tau_8}.\notag
\end{align}
For the sake of convenience, we define a new parameter $\tau_9$ as follows:
\begin{equation*}
\tau_9=
\begin{cases}
\tau_6, & \text{ if  }\frac{6-2p}{6-p} + \frac{6-3q}{6-q} < 1, \\
\tau_8, & \text{ if  } \frac{6-2p}{6-p} + \frac{6-3q}{6-q}\geq1\text{ and }\gamma=0,\\
\frac{1}{2}, & \text{ if  } \frac{6-2p}{6-p} + \frac{6-3q}{6-q}\geq1\text{ and }\gamma>0.
\end{cases}
\end{equation*}
Then putting \eqref{ine4.38}, \eqref{ine4.39} and \eqref{ine4.40} together, we obtain
\begin{align}\label{ine4.41}
K_{711}\leq&C\left[R\ln RE'(R)\right]^{\tau_9}.
\end{align}


Next, we deal with the terms $K_{712}$ and $K_{713}$. For $\frac{3}{2}<p<3$ and $R>3$, we claim
 $$
 R^{\frac{p}{6-p}\alpha+\frac{2q}{6-q}\beta-1+\frac{6-3q}{6-q}\left[\beta-\left(\frac{3}{q}-\frac{1}{2}\right)\right]}(\ln R)^{\frac{p}{6-p}\lambda+\mu}\leq C(\ln R)^{\frac{3-p}{6-p}+\frac{1}{2}}.
 $$
Indeed, when $q=2$, using Assumption \ref{a1.4}(i), we have
 $$
 R^{\frac{p}{6-p}\alpha+\frac{2q}{6-q}\beta-1+\frac{6-3q}{6-q}\left[\beta-\left(\frac{3}{q}-\frac{1}{2}\right)\right]}(\ln R)^{\frac{p}{6-p}\lambda+\mu}\leq C(\ln R)^{\frac{p}{6-p}\lambda+\mu}\leq C(\ln R)^{\frac{3-p}{6-p}+\frac{1}{2}}.
 $$
 When $1\leq q<2$ and $\beta<\frac{3}{q}-\frac{1}{2}$, we find
 $$
 R^{\frac{p}{6-p}\alpha+\frac{2q}{6-q}\beta-1+\frac{6-3q}{6-q}\left[\beta-\left(\frac{3}{q}-\frac{1}{2}\right)\right]}(\ln R)^{\frac{p}{6-p}\lambda+\mu}\leq C\leq C(\ln R)^{\frac{3-p}{6-p}+\frac{1}{2}}.
 $$
 When $1\leq q<2$ and $\beta=\frac{3}{q}-\frac{1}{2}$, using Assumption \ref{a1.4}(iii), we get
 $$
 R^{\frac{p}{6-p}\alpha+\frac{2q}{6-q}\beta-1+\frac{6-3q}{6-q}\left[\beta-\left(\frac{3}{q}-\frac{1}{2}\right)\right]}(\ln R)^{\frac{p}{6-p}\lambda+\mu}\leq C(\ln R)^{\frac{p}{6-p}\lambda+2\mu}\leq C(\ln R)^{\frac{3-p}{6-p}+\frac{1}{2}}.
 $$
Employing \eqref{ine4.15} and \eqref{ine4.2}, we obtain
\begin{align}\label{ine4.42}
	K_{712}\leq&C R^{\frac{p}{6-p}\alpha+\frac{2q}{6-q}\beta-1+\frac{6-3q}{6-q}\left[\beta-\left(\frac{3}{q}-\frac{1}{2}\right)\right]}(\ln R)^{\frac{p}{6-p}\lambda+\mu}\left[RE'(R)\right]^{\frac{3-p}{6-p}+\frac{1}{2}}\notag\\
  \leq&C\left[R\ln RE'(R)\right]^{\frac{3-p}{6-p}+\frac{1}{2}}.
\end{align}
Employing \eqref{ine4.15} and \eqref{ine4.14}, we obtain
\begin{align*}
K_{713} \leq&C R^{-1} \|u\|_{L^p(A_R)}^{\frac{p}{6-p}}\|v\|_{L^q(A_R)}^{\frac{2q}{6-q}}\left(R^{\frac{1}{2}-\frac{3}{p}} \|u\|_{L^p(A_R)}\right)^{\frac{6-2p}{6-p}}\left[f(2R)\right]^{\frac{6-2q}{6-q}}\\
\leq&CR^{\frac{6-2p}{6-p}\left[\alpha-\left(\frac{3}{p}-\frac{1}{2}\right)\right]}(\ln R)^{\lambda+\frac{2q}{6-q}\mu+\frac{6-2q}{6-q}\tau_5},
\end{align*}
which implies $\lim\limits_{R\rightarrow+\infty}K_{713}=0$.
Consequently, there exists a constant $R_7>R_2$ such that
\begin{align}\label{ine4.43}
K_{713}\leq\frac{1}{16}E(R_2)\leq \frac{1}{16}E(R),\;\forall R\geq R_7.
\end{align}
Combining \eqref{ine4.36a}, \eqref{ine4.37}, \eqref{ine4.41}, \eqref{ine4.42} and \eqref{ine4.43}, we have
\begin{align}\label{ine4.44}
K_{71}\leq&\frac{1}{16}E(R)+C\left[R\ln RE'(R)\right]^{\tau_9}+C\left[R\ln RE'(R)\right]^{\frac{3-p}{6-p}+\frac{1}{2}}\notag\\
&+C\left[R\ln RE'(R)\right]^{\frac{6-3q}{12-2q}+\frac{1}{2}}+C\left[R\ln R E'(R)\right]^{\frac{1}{2}},\;\forall R\geq R_7.
\end{align}



We then turn to handle $K_{72}$.
When $\beta<\frac{3}{q}-\frac{1}{2}$, by the interpolation inequality, \eqref{ine4.15} and \eqref{ine4.14}, we obtain
\begin{align*}
	K_{72} \leq& C R^{-\frac{3}{q}-\frac{1}{2}} \|u\|_{L^p(A_R)}^{\frac{p}{6-p}}\|v\|_{L^q(A_R)}^{\frac{2q}{6-q}} \|u\|_{L^6(A_R)}^{\frac{6-2p}{6-p}}\|v\|_{L^6(A_R)}^{\frac{6-3q}{6-q}}  \|v\|_{L^q(A_R)}\\
 \leq&C R^{-\frac{3}{q}-\frac{1}{2}}\|u\|_{L^p(A_R)}^{\frac{p}{6-p}}\|v\|_{L^q(A_R)}^{\frac{2q}{6-q}}\left[f(2R)\right]^{\frac{3-p}{6-p}+\frac{6-3q}{12-2q}}\|v\|_{L^q(A_R)}\\
 \leq&C R^{\frac{p}{6-p}\alpha+\frac{2q}{6-q}\beta-1+\beta-\left(\frac{3}{q}-\frac{1}{2}\right)}(\ln R)^{\frac{p}{6-p}\lambda+\frac{6+q}{6-q}\mu+\left(\frac{3-p}{6-p}+\frac{6-3q}{12-2q}\right)\tau_5}\\
 \leq&C R^{\frac{1}{2}\left[\beta-\left(\frac{3}{q}-\frac{1}{2}\right)\right]},
\end{align*}
which forces $\lim\limits_{R\rightarrow+\infty}K_{72}=0$. Hence, there exists a constant $R_8>R_2$ such that
\begin{align}\label{ine4.45}
K_{72}\leq\frac{1}{16}E(R_2)\leq \frac{1}{16}E(R),\;\forall R\geq R_8.
\end{align}


When $\beta=\frac{3}{q}-\frac{1}{2}$, we adopt a different strategy to handle $K_{72}$. It is noted that when $\beta=\frac{3}{q}-\frac{1}{2}$, we require $p\neq3$.
By the interpolation inequality, \eqref{ine4.19} and \eqref{ine4.30}, we obtain
\begin{align}\label{ine4.46}
	K_{72} \leq& CR^{-\frac{3}{q}-\frac{1}{2}}\|u\|_{L^p(A_R)}^{\frac{p}{6-p}}\|v\|_{L^q(A_R)}^{\frac{2q}{6-q}} \|u\|_{L^6(A_R)}^{\frac{6-2p}{6-p}}\|v\|_{L^6(A_R)}^{\frac{6-3q}{6-q}} \|v\|_{L^q(A_R)}\notag\\
	\leq &CR^{-\frac{3}{q}-\frac{1}{2}}\|u\|_{L^p(A_R)}^{\frac{p}{6-p}}\|v\|_{L^q(A_R)}^{\frac{2q}{6-q}}\|\nabla u\|_{L^2(A_R)}^{\frac{6-2p}{6-p}}\|\nabla v\|_{L^2(A_R)}^{\frac{6-3q}{6-q}}\|v\|_{L^q(A_R)}\notag\\
&+CR^{-\frac{3}{q}-\frac{1}{2}}\|u\|_{L^p(A_R)}^{\frac{p}{6-p}}\|v\|_{L^q(A_R)}^{\frac{2q}{6-q}}\|\nabla u\|_{L^2(A_R)}^{\frac{6-2p}{6-p}}\left(R^{\frac{1}{2}-\frac{3}{q}} \|v\|_{L^q(A_R)}\right)^{\frac{6-3q}{6-q}}\|v\|_{L^q(A_R)}\\
&+CR^{-\frac{3}{q}-\frac{1}{2}}\|u\|_{L^p(A_R)}^{\frac{p}{6-p}}\|v\|_{L^q(A_R)}^{\frac{2q}{6-q}}\left(R^{\frac{1}{2}-\frac{3}{p}} \|u\|_{L^p(A_R)}\right)^{\frac{6-2p}{6-p}} \|v\|_{L^6(A_R)}^{\frac{6-3q}{6-q}}\|v\|_{L^q(A_R)}\notag\\
=&:K_{721} + K_{722} + K_{723}.\notag
\end{align}
Using \eqref{ine4.15}, \eqref{ine4.2} and Assumption \ref{a1.4}(iii), we obtain
\begin{align}\label{ine4.47}
K_{721}\leq&CR^{\frac{p}{6-p}\alpha+\frac{2q}{6-q}\beta-1+\beta-\left(\frac{3}{q}-\frac{1}{2}\right)}(\ln R)^{\frac{p}{6-p}\lambda+\frac{6+q}{6-q}\mu} \left[RE'(R)\right]^{\frac{3-p}{6-p}+\frac{6-3q}{12-2q}}\notag\\
\leq&C(\ln R)^{\frac{p}{6-p}\lambda+2\mu} \left[RE'(R)\right]^{\frac{3-p}{6-p}+\frac{6-3q}{12-2q}}\notag\\
\leq&C(\ln R)^\frac{3-p}{6-p} \left[RE'(R)\right]^{\frac{3-p}{6-p}+\frac{6-3q}{12-2q}}\\
\leq&C\left[R\ln RE'(R)\right]^{\frac{3-p}{6-p}+\frac{6-3q}{12-2q}},\notag
\end{align}
and
\begin{align}\label{ine4.48}
K_{722}\leq&CR^{\frac{p}{6-p}\alpha+\frac{2q}{6-q}\beta-1+\frac{12-4q}{6-q}\left[\beta-\left(\frac{3}{q}-\frac{1}{2}\right)\right]}(\ln R)^{\frac{p}{6-p}\lambda+2\mu} \left[RE'(R)\right]^{\frac{3-p}{6-p}}\notag\\
\leq&C\left[R\ln RE'(R)\right]^{\frac{3-p}{6-p}}.
\end{align}
Using \eqref{ine4.15} and \eqref{ine4.14}, we obtain
\begin{align*}
K_{723}\leq&CR^{\frac{6-2p}{6-p}\left[\alpha-\left(\frac{3}{p}-\frac{1}{2}\right)\right]+\frac{p}{6-p}\alpha+\frac{2q}{6-q}\beta-1+\beta-\left(\frac{3}{q}-\frac{1}{2}\right)}(\ln R)^{\lambda+\frac{6+q}{6-q}\mu} \left[f(2R)\right]^{\frac{6-3q}{12-2q}}\\
\leq&CR^{\frac{6-2p}{6-p}\left[\alpha-\left(\frac{3}{p}-\frac{1}{2}\right)\right]}(\ln R)^{\lambda+\frac{6+q}{6-q}\mu+\frac{6-3q}{12-2q}\tau_5},
\end{align*}
which implies $\lim\limits_{R\rightarrow+\infty}K_{723}=0$.
Thus, there exists a constant $R_9>R_2$ such that
\begin{align}\label{ine4.49}
	K_{723}\leq\frac{1}{16}E(R_2)\leq \frac{1}{16}E(R),\;\forall R\geq R_9.
\end{align}
We define a constant $R_{10}$ by
\begin{equation*}
R_{10}=
\begin{cases}
R_8, & \text{ if  }\;0\leq\beta<\frac{3}{q}-\frac{1}{2}, \\
R_9, & \text{ if  }\;\beta=\frac{3}{q}-\frac{1}{2}.
\end{cases}
\end{equation*}
Collecting \eqref{ine4.45}, \eqref{ine4.46}, \eqref{ine4.47}, \eqref{ine4.48} and \eqref{ine4.49},
we conclude that for any $R\geq R_{10}$, it holds that
\begin{align}\label{ine4.50}
K_{72}\leq \frac{1}{16}E(R)+C\chi(p)\left(\left[R\ln RE'(R)\right]^{\frac{3-p}{6-p}+\frac{6-3q}{12-2q}}+\left[R\ln RE'(R)\right]^{\frac{3-p}{6-p}}\right).
\end{align}

Using \eqref{ine4.8}, \eqref{ine4.16} and \eqref{ine4.17}, we derive
\begin{align}\label{ine4.51}
	K_{73} &\leq C\|\nabla v\|_{L^2 (A_R)}^2+CR^{\frac{1}{2}-\frac{3}{q}}\|v\|_{L^q(A_R)}\|\nabla v\|_{L^2 (A_R)}\notag\\
           &\leq C[R\ln RE'(R)]^\frac{\tau_5}{1+\tau_5}+C[R\ln RE'(R)]^\frac{1}{2}.
\end{align}
We denote $R_{11}=\max\{R_7,R_{10}\}$. Combining \eqref{ine4.44}, \eqref{ine4.50} and \eqref{ine4.51}, we get \eqref{ine4.35}.
\end{proof}

\end{document}
